\section{Introduction}
Communication edges are useful only through their effects on downstream
answers. In a team of small language models, those effects depend on which
other messages remain available. A topology edit can therefore invalidate
a previously measured edge credit. The practical question is not simply
whether credit changes, but whether a limited feedback budget can identify
changes that alter useful decisions and recover more service value than
the feedback costs.

Three claims must be distinguished. A predictor may estimate conditional
credit more accurately; a graph change may measurably affect credit; and a
policy may profit from buying new measurements. None entails the next.
For sparse credit outcomes, predicting zero can minimize aggregate squared
error while missing decision-relevant events. Conversely, even a clear
credit change may already be handled by a free structural rule, leaving
little value for paid revalidation.

We investigate these distinctions in a bounded research cycle rather than
presenting an unqualified new optimization method. The original credit
prediction study failed its gate. Diagnostic controls identified a constant
response artifact from intercept regularization, yet a fixed corrected
predictor did not improve exposed editing decisions. Subsequent controlled
path experiments detected topology-sensitive credit while their nominated
graphs underperformed inexpensive system baselines. These observations
motivated a final candidate that explicitly prices initialization,
counterfactual feedback, and graph execution.

The final hypothesis was that graph-dependent execution costs and already
purchased request reuse could make decision-focused revalidation outperform
both ordinary acquisition and free controls. We fixed the predictor, data
partition, candidate, cost weight, baselines, and admission criteria before
revealing the new cohort's scores. All methods share exact omission of
nodes outside the answer node's ancestor set. This execution principle is
not attributed to the proposed acquisition policy.

Our evidence is negative: the final nominee fails four of six conjunctive
development criteria. The bounded contribution is a reproducible failure
analysis and evaluation design, including paid-information accounting,
strong inexpensive controls, and separate numerical reconstruction. We do
not claim a successful new method, external state-of-the-art performance,
or a general impossibility result. Appendix~\ref{app:history} retains the
historical decisions instead of selecting only favorable experiments.
